\chapter{Introduction}\label{ch:intro}

\section{What \RITMO{} is}

The Atari XL/XE computers make their sound with the POKEY chip: four independent voices
(sixteen when two chips are used, as in many stereo songs), each with a frequency divider, a
choice of noise and tone distortions and a volume. A \emph{tracker} is the kind of music editor
that was born on the Amiga: the music is a grid of numbers, one column for each voice and one
row for each beat, and the notes are typed with the computer keyboard as on a piano.

\RITMO{} is such a tracker for the POKEY. A song made with it can be played on a real Atari: the
program produces the same data the player routines of the Atari read, and exports them in
several formats (a SAP file for the Atari players, an executable \texttt{.xex}, assembler source,
the compact RMT module for your own program). The sound you hear while composing comes from an
emulation of the Atari's 6502 processor and of the POKEY chip built into the program, running
the same player routine, so it is what the Atari will play.

\shot{window-tracks}{The main window of \RITMO{} with a song loaded.}{0.98}

\section{History}

\RITMO{} is a fork of the \emph{RASTER Music Tracker}, \RMT{}. The first version, 1.0, was
written by Radek \v{S}t\v{e}rba (Raster/C.P.U.) in 2002; his last version is 1.28, of 2009. It was
a small revolution for the Atari musicians: the first tracker for the POKEY with a
graphical envelope editor, and the one nearly everybody used. Radek \v{S}t\v{e}rba has since
passed away; this manual is, among other things, a way to remember his work. R.I.P.

Vin Samuel (\emph{VinsCool}) continued the program from 2021 to 2024, up to the version 1.34: a
new player routine with more tuning tables, a tuning calculator, an automatic high pass filter,
the SAP and XEX export and many corrections. Peter Dell (\emph{JAC!}) carries on the original
project from 2024 (version 1.35 and the planned version 2) and made it possible for this fork
to exist: everything described here is built on his work.

The \RITMO{} fork keeps the code of the tracker, the file formats and the player routines of
\RMT{} 1.35, and replaces what tied it to Windows. It is renamed so that it cannot be confused with the original
\RMT{}, and so that the Windows-only (MFC) parts of the code could be removed.

\section{What is different from \RMT{}}

\begin{itemize}
  \item The program is written with Qt: it runs natively on Linux, on Windows and on macOS (Intel and
        Apple Silicon). The look of the tracker screen is the original one, pixel for pixel.
  \item The 6502 and the POKEY emulations are built in: the \texttt{sa\_c6502.dll} and
        \texttt{apokeysnd.dll} libraries of the Windows versions are not needed any more.
  \item The sound goes out through PortAudio, and the MIDI input comes in through RtMidi.
  \item The configuration is kept per user in the settings of the system, not in an \texttt{.ini}
        file next to the program. The settings of \RMT{} are taken over at the first start.
  \item Songs, instruments and tracks are the same files as in \RMT{} 1.3x: you can move your work
        between the programs. The files keep their extensions, \texttt{.rmt} and the others.
\end{itemize}

\section{Files used by \RITMO{}}

\begin{center}
\begin{tabular}{@{}p{3.6cm}p{10.4cm}@{}}
\toprule
Extension & Content \\ \midrule
\texttt{.rmt} & an RMT module: the song, its tracks and its instruments (and, if the song was saved
                with \menu{Save}, the names of the song and of the instruments) \\
\texttt{.txt} & the same song as a text file, which can be read and edited with any editor; a single track is saved as a text file too \\
\texttt{.rmw} & a song work file: an RMT module with the state of the editing \\
\texttt{.rti} & a single instrument \\
\texttt{.sap}, \texttt{.sapr}, \texttt{.xex}, \texttt{.asm}, \texttt{.wav}, \texttt{.lzss} & the exported formats (chapter~\ref{ch:importexport}) \\
\texttt{.mod}, \texttt{.tmc}, \texttt{.tm8} & the formats that can be imported \\
\texttt{.rmtscript} & a script (chapter~\ref{ch:scripting}) \\
\bottomrule
\end{tabular}
\end{center}
